\section{The Gaussian quadratic form in a sector}
\label{sec:sector-gram}
The quadratic reference must remain nondegenerate after the large-field
bonds have been frozen. A lower bound for a flat Hessian alone would
not establish this. We prove the lower bound for the actual transported,
selected rows, using a separate local gauge around each active bond.
All norms in the quadratic inequalities below are physical counting
norms: squared colour norms are summed over bonds without dividing
by the volume.

Fix the tree, seed set, coarse field and original predictor tags of
Section~\ref{sec:packet-annulus}. Let $\mathcal H$ be the active
free bonds and let $P_{\mathcal H}$ denote restriction to them,
extended by zero on all other bonds. The seed store's coarse-flag
component equals the actual coarse flags; it may have other seed
components. An active bond has both coarse endpoints outside
the $10B$ seed neighbourhood. A row is selected when its anchor
is outside the $6B$ seed neighbourhood. These are fixed selections
for this quadratic form.

Let $q$ be the number of signed channels and write
$\lambda=q^{-1/2}$. For a channel $s$ and a fine or predictor
row $a$, set
\[
 D_{s,a}u={\bf1}_{\{(s,a)\ \text{selected}\}}
             A^V_{s,a}P_{\mathcal H}u.
\]
The map $A^V_{s,a}$ is the flat linear cut row transported using $V$,
including the signed displacement of input colours. It is not, in general,
the derivative at the replica of $V$ of the nonlinear curvature or predictor. In
orthonormal colour coordinates, define $M$ by
\[
 (Mu)_{s,a}=\lambda D_{s,a}u.
\]
Here and below a colour array and its scalar coordinate vector are
identified. Define the completed precision matrix
\begin{equation}\label{eq:sector-precision}
 A=M^*M+P_{\mathcal H^c}.
\end{equation}
Every column of $M$ indexed by an inactive bond is zero.
Consequently $A$ acts as the identity on the inactive coordinates
and has no cross terms between them and the active coordinates.
The completion supplies no energy on the active subspace.

If $b$ is the vector of selected affine rows at zero fluctuation,
including the same factor $\lambda$, linearity of the transported
rows gives
\[
 Q_{\rm quad}(u)=\tfrac12\|b+Mu\|_2^2.
\]
This is exactly the orientation average of the extracted squares:
$\lambda^2=1/q$ supplies its normalization. In particular,
\[
 Q_{\rm quad}(v+u)-Q_{\rm quad}(v)
 =\langle u,M^*b+M^*Mv\rangle
       +\tfrac12\langle u,M^*Mu\rangle.
\]
The same fine field, tags and orientation choices are used in every
term. We have not averaged independently chosen Gaussian matrices.

\subsection{A local linear inequality from the tree comparison}
Let $w$ be a fine colour array which vanishes on the fixed tree.
The tree comparison of Section~\ref{sec:seed-tree-patch}, with
coarse configuration equal to the identity, implies
\begin{equation}\label{eq:sector-local-linear}
 \|w_e\|\le K\sum_{z\in\mathcal B}\sum_\alpha
       \|(d_fw)_{\overline z+L o,\alpha}\|
           +\|(\mathcal Qw)_{o,a}\|,
 \qquad o=\operatorname{parent}(e),
\end{equation}
where $a$ is the axis of $e$, $\mathcal B$ is the fixed macro box,
and $\mathcal Q$ is the derivative of the actual predictor.
The box sum may count periodic images more than once; this only
enlarges its nonnegative right side.

Here is the passage to the linear inequality. Apply the nonlinear
comparison to $U_e(\varepsilon)=\exp(\varepsilon w_e)$.
It remains in the same tree slice since $w$ vanishes on the tree.
For sufficiently small $\varepsilon$ all required frame
likelihoods are positive: at zero every frame is in the minimizing
branch, and the finitely many conditions are open. Divide the
comparison by $|\varepsilon|$ and let $\varepsilon\to0$ through
nonzero real values. If $g(0)=1$ and its matrix derivative is the
colour matrix of $v$, then
\[
 \lim_{\varepsilon\to0}
 \frac{r(g(\varepsilon))}{|\varepsilon|}=\|v\|.
\]
This follows directly from differentiability and the physical
normalization of the chord. The derivatives of the fine links,
plaquettes and predictor are respectively $w$, $d_fw$ and
$\mathcal Qw$, the last by Section~\ref{sec:predictor-branch}.
Passing through the finite box sum proves
\eqref{eq:sector-local-linear}. The limiting neighbourhood can
depend on the finite torus; the resulting constant $K$, inherited
from the nonlinear tree inequality, does not.

\subsection{A lower bound for the actual selected rows}
For an anchor $o$, define
\[
 \begin{split}
 J_o(u)^2&=\sum_{\operatorname{dist}(o,\operatorname{parent}(e))
                         \le3B}\|(P_{\mathcal H}u)_e\|^2,\\
 R_o(u)^2&=\sum_{\operatorname{dist}(o,\operatorname{anchor}(s,a))
                         \le R_{\rm macro}+2}\|D_{s,a}u\|^2.
 \end{split}
\]
The square roots are nonnegative. An active bond at $o$ is more
than $10B$ from the seeds. Since the macro radius is much smaller
than $B$, every row with anchor within $R_{\rm macro}+2$ of $o$
is selected in the identity signed channel. Also the coarse patch
of radius $2B+1$ is clean, because the coarse flags belong to the
seeds. The local axial construction gives a gauge $g$ with
\[
 r((g\cdot V)_e)\le h_0(x)
 \quad\hbox{on radius }2B,
\]
where $h_0\ge0$ has upper exponent $-1/2$.

Apply \eqref{eq:sector-local-linear} to
$w=\operatorname{Ad}(g)P_{\mathcal H}u$. This array still
vanishes on the tree. For every required row, the transport estimate
\eqref{eq:annulus-transport-error} and gauge covariance give
\begin{equation}\label{eq:sector-flat-row-control}
 \|(A_Rw)_a\|\le R_o(u)+\varepsilon(x)J_o(u),
 \qquad \varepsilon(x)=C\,2D_{\rm col}(B+1)h_0(x).
\end{equation}
Indeed the actual row in the chosen gauge has the same norm as
its selected original row and is at most $R_o(u)$. Every input
used in the comparison lies within $3B$ of $o$, so its norm is
at most $J_o(u)$. All comparison routes lie inside the radius
$2B$ patch. The identity signed channel suffices for this lower
bound; the sum defining $R_o$ may include all the other channels.

The direct fine component of $A_Rw$ is $\sqrt{a_f}\,d_fw$,
and its predictor component is $-\sqrt p\,\mathcal Qw$.
Writing $N_{\rm macro}=|\mathcal B|N_{\rm plaq}$, insert
\eqref{eq:sector-flat-row-control} into
\eqref{eq:sector-local-linear}. We obtain
\begin{equation}\label{eq:sector-bond-control}
 \|(P_{\mathcal H}u)_e\|
 \le ZR_o(u)+E(x)J_o(u),
 \quad Z=K N_{\rm macro}/\sqrt{a_f}+1/\sqrt p,
 \quad E=Z\varepsilon.
\end{equation}
This also holds on inactive bonds because its left side is zero.
$Z$ is fixed after $L$; $E$ has upper exponent $-1/2$.

There are $L^4N_{\rm axis}$ fine positive bonds per coarse site.
Let
\[
 S_L(d)=(2d+1)^4L^4N_{\rm axis}.
\]
For every nonnegative family $f_a$ anchored on coarse sites,
changing the order of the two finite sums shows
\[
 \sum_e\sum_{\operatorname{dist}(\operatorname{parent}(e),
                                    \operatorname{anchor}(a))\le d} f_a
 \le S_L(d)\sum_a f_a.
\]
For fixed $a$, at most $(2d+1)^4$ coarse sites can occur, each
with at most $L^4N_{\rm axis}$ bonds. This remains valid when
the neighbourhood wraps around the torus: wrapping reduces the
number of distinct sites.

Square \eqref{eq:sector-bond-control}, using
$(r+s)^2\le2r^2+2s^2$, and sum over bonds. The counting bound gives
\[
 \|P_{\mathcal H}u\|_2^2
 \le2Z^2S_L(R_{\rm macro}+2)\sum_{s,a}\|D_{s,a}u\|^2
       +2E(x)^2S_L(3B)\|P_{\mathcal H}u\|_2^2.
\]
The last coefficient has upper exponent $-1$ and is eventually
at most $1/2$. Move it to the left to get
\[
 \|P_{\mathcal H}u\|_2^2
 \le a_*\sum_{s,a}\|D_{s,a}u\|^2,
 \qquad a_*=1+4Z^2S_L(R_{\rm macro}+2).
\]
The added one ensures a strictly positive constant in every case.
Since $\|Mu\|_2^2=\lambda^2\sum\|D_{s,a}u\|^2$, the
orthogonal active/inactive decomposition proves
\begin{equation}\label{eq:sector-gram-floor}
 \langle u,Au\rangle\ge\kappa\|u\|_2^2,
 \qquad \kappa=\min\{\lambda^2/a_*,1\}>0.
\end{equation}
This includes an empty active set, in which case $A=1$.
No globally small gauge for $V$ has been constructed or assumed.

\subsection{Upper bounds with spatial weights}
We also need a bound independent of the period on the matrix
entries and their weighted sums. The cut row $A_R$ vanishes on
an impulse at a bond farther than $R+\rho+4$ from its anchor.
Together with \eqref{eq:annulus-linear-rows}, this implies
\[
 e^{\sigma d}\|(A_R\delta_ev)_a\|
 \le2C e^{\sigma(\rho+6)}\|v\|,
 \qquad d=\operatorname{dist}(\operatorname{anchor}(a),
                                      \operatorname{parent}(e)).
\]
To see this, if $d\le\rho+5$ use the uniform cut-row bound.
Otherwise take $r=d-(\rho+5)$, so that $A_r\delta_ev=0$.
If $R<r$ locality gives zero as well. For $R\ge r$, compare
both $A_R$ and $A_r$ with $A_\infty$; the two tails are at most
$C\|v\|e^{-\sigma r}$ each. This proves the stated bound.

Signed displacement and parallel transport send a single impulse
to a single impulse with unchanged colour norm. The coarse anchor
may move by one under reversal of a bond. Thus the preceding
constant becomes $2Ce^{\sigma(\rho+7)}$ for the actual oriented
row. Selection and restriction to active bonds only set entries
to zero. In scalar coordinates,
\[
 e^{\sigma\operatorname{dist}(i,j)}|M_{ij}|
 \le b_*=\lambda\,2Ce^{\sigma(\rho+7)}.
\]
The anchors of input coordinates have multiplicity at most
$m=L^4N_{\rm axis}q_{\rm col}$, where $q_{\rm col}$ is the
colour dimension. Output-coordinate anchors have multiplicity
at most
\[
 n=q\,(2L^4N_{\rm plaq}+N_{\rm axis})q_{\rm col}.
\]
For $\eta>0$ put
$P_0(\eta)=\sum_{k\ge0}(2k+1)^4e^{-\eta k}<\infty$.
Grouping coordinates by their distance from an anchor gives,
for $0\le r\le\sigma/2$,
\[
 \sup_i\sum_j e^{r\operatorname{dist}(i,j)}|M_{ij}|
 \le b_*mP_0(\sigma/2),\qquad
 \sup_j\sum_i e^{r\operatorname{dist}(i,j)}|M_{ij}|
 \le b_*nP_0(\sigma/2).
\]
Choose one constant $C_*$ larger than these two bounds and one.
The triangle inequality for anchor distances proves the product
rule for such weighted row and column sums. Therefore
\begin{equation}\label{eq:sector-gram-upper}
 \sup_i\sum_j e^{r\operatorname{dist}(i,j)}|A_{ij}|
 \le C_*^2+1.
\end{equation}
The extra one is the inactive diagonal. Symmetry gives the same
column bound. The usual finite Cauchy--Schwarz argument gives
$\|A\|_{2\to2}\le C_*^2+1$: bound
$|(Au)_i|^2$ by the absolute row sum times
$\sum_j|A_{ij}||u_j|^2$ and then sum over $i$.
Also $M_{ij}=0$ beyond distance $B/4$ by its stencil, so
$A_{ij}=0$ beyond $B/2$.

Combining \eqref{eq:sector-gram-floor} and
\eqref{eq:sector-gram-upper} proves a strictly positive, uniform
spectral interval for the actual completed precision matrix.
In particular its inverse exists and satisfies
$0<A^{-1}\le\kappa^{-1}1$. For $\beta>0$, the centered finite-dimensional
Gaussian with covariance $\beta^{-1}A^{-1}$ is consequently
well defined, including when there are no active coordinates.
Restricting to the active subspace gives the physical Gaussian;
the inactive identity coordinates are an algebraic completion.
Spatial locality of the inverse, comparison of local Gaussian
normalizations, and control of the affine mean require the
additional estimates developed at their respective steps.
